% Options for packages loaded elsewhere
% Options for packages loaded elsewhere
\PassOptionsToPackage{unicode}{hyperref}
\PassOptionsToPackage{hyphens}{url}
\PassOptionsToPackage{dvipsnames,svgnames,x11names}{xcolor}
%
\documentclass[
  12pt,
]{article}
\usepackage{xcolor}
\usepackage[top=25mm,bottom=25mm,left=25mm,right=25mm]{geometry}
\usepackage{amsmath,amssymb}
\setcounter{secnumdepth}{5}
\usepackage{iftex}
\ifPDFTeX
  \usepackage[T1]{fontenc}
  \usepackage[utf8]{inputenc}
  \usepackage{textcomp} % provide euro and other symbols
\else % if luatex or xetex
  \usepackage{unicode-math} % this also loads fontspec
  \defaultfontfeatures{Scale=MatchLowercase}
  \defaultfontfeatures[\rmfamily]{Ligatures=TeX,Scale=1}
\fi
\usepackage{lmodern}
\ifPDFTeX\else
  % xetex/luatex font selection
\fi
% Use upquote if available, for straight quotes in verbatim environments
\IfFileExists{upquote.sty}{\usepackage{upquote}}{}
\IfFileExists{microtype.sty}{% use microtype if available
  \usepackage[]{microtype}
  \UseMicrotypeSet[protrusion]{basicmath} % disable protrusion for tt fonts
}{}
\usepackage{setspace}
\makeatletter
\@ifundefined{KOMAClassName}{% if non-KOMA class
  \IfFileExists{parskip.sty}{%
    \usepackage{parskip}
  }{% else
    \setlength{\parindent}{0pt}
    \setlength{\parskip}{6pt plus 2pt minus 1pt}}
}{% if KOMA class
  \KOMAoptions{parskip=half}}
\makeatother
% Make \paragraph and \subparagraph free-standing
\makeatletter
\ifx\paragraph\undefined\else
  \let\oldparagraph\paragraph
  \renewcommand{\paragraph}{
    \@ifstar
      \xxxParagraphStar
      \xxxParagraphNoStar
  }
  \newcommand{\xxxParagraphStar}[1]{\oldparagraph*{#1}\mbox{}}
  \newcommand{\xxxParagraphNoStar}[1]{\oldparagraph{#1}\mbox{}}
\fi
\ifx\subparagraph\undefined\else
  \let\oldsubparagraph\subparagraph
  \renewcommand{\subparagraph}{
    \@ifstar
      \xxxSubParagraphStar
      \xxxSubParagraphNoStar
  }
  \newcommand{\xxxSubParagraphStar}[1]{\oldsubparagraph*{#1}\mbox{}}
  \newcommand{\xxxSubParagraphNoStar}[1]{\oldsubparagraph{#1}\mbox{}}
\fi
\makeatother


\usepackage{longtable,booktabs,array}
\usepackage{calc} % for calculating minipage widths
% Correct order of tables after \paragraph or \subparagraph
\usepackage{etoolbox}
\makeatletter
\patchcmd\longtable{\par}{\if@noskipsec\mbox{}\fi\par}{}{}
\makeatother
% Allow footnotes in longtable head/foot
\IfFileExists{footnotehyper.sty}{\usepackage{footnotehyper}}{\usepackage{footnote}}
\makesavenoteenv{longtable}
\usepackage{graphicx}
\makeatletter
\newsavebox\pandoc@box
\newcommand*\pandocbounded[1]{% scales image to fit in text height/width
  \sbox\pandoc@box{#1}%
  \Gscale@div\@tempa{\textheight}{\dimexpr\ht\pandoc@box+\dp\pandoc@box\relax}%
  \Gscale@div\@tempb{\linewidth}{\wd\pandoc@box}%
  \ifdim\@tempb\p@<\@tempa\p@\let\@tempa\@tempb\fi% select the smaller of both
  \ifdim\@tempa\p@<\p@\scalebox{\@tempa}{\usebox\pandoc@box}%
  \else\usebox{\pandoc@box}%
  \fi%
}
% Set default figure placement to htbp
\def\fps@figure{htbp}
\makeatother





\setlength{\emergencystretch}{3em} % prevent overfull lines

\providecommand{\tightlist}{%
  \setlength{\itemsep}{0pt}\setlength{\parskip}{0pt}}



 


% Line numbering for co-author review.
% Remove \linenumbers before final journal submission.
\usepackage{lineno}
\linenumbers
\modulolinenumbers[5]

% Paragraph spacing
\setlength{\parindent}{0em}
\setlength{\parskip}{0.8em}
\makeatletter
\@ifpackageloaded{caption}{}{\usepackage{caption}}
\AtBeginDocument{%
\ifdefined\contentsname
  \renewcommand*\contentsname{Table of contents}
\else
  \newcommand\contentsname{Table of contents}
\fi
\ifdefined\listfigurename
  \renewcommand*\listfigurename{List of Figures}
\else
  \newcommand\listfigurename{List of Figures}
\fi
\ifdefined\listtablename
  \renewcommand*\listtablename{List of Tables}
\else
  \newcommand\listtablename{List of Tables}
\fi
\ifdefined\figurename
  \renewcommand*\figurename{Figure}
\else
  \newcommand\figurename{Figure}
\fi
\ifdefined\tablename
  \renewcommand*\tablename{Table}
\else
  \newcommand\tablename{Table}
\fi
}
\@ifpackageloaded{float}{}{\usepackage{float}}
\floatstyle{ruled}
\@ifundefined{c@chapter}{\newfloat{codelisting}{h}{lop}}{\newfloat{codelisting}{h}{lop}[chapter]}
\floatname{codelisting}{Listing}
\newcommand*\listoflistings{\listof{codelisting}{List of Listings}}
\makeatother
\makeatletter
\makeatother
\makeatletter
\@ifpackageloaded{caption}{}{\usepackage{caption}}
\@ifpackageloaded{subcaption}{}{\usepackage{subcaption}}
\makeatother
\usepackage{bookmark}
\IfFileExists{xurl.sty}{\usepackage{xurl}}{} % add URL line breaks if available
\urlstyle{same}
\hypersetup{
  pdftitle={TODO: manuscript title},
  pdfauthor={Ondřej Mottl; TODO Co-author 2; TODO Co-author 3},
  pdfkeywords={vegetation co-occurrence, palaeoecology, joint species
distribution model, HMSC, Holocene},
  colorlinks=true,
  linkcolor={blue},
  filecolor={Maroon},
  citecolor={Blue},
  urlcolor={Blue},
  pdfcreator={LaTeX via pandoc}}


\title{TODO: manuscript title}
\author{Ondřej Mottl \and TODO Co-author 2 \and TODO Co-author 3}
\date{April 29, 2026}
\begin{document}
\maketitle


\setstretch{2}
\subsection{Introduction}\label{introduction}

TODO: Write introduction.

\subsection{Methods}\label{methods}

\textsubscript{Source:
\href{https://OndrejMottl.github.io/BIODYNAMICS-vegetation-cooccurrence/manuscript.qmd.html}{Article
Notebook}}

\subsubsection{Vegetation and palaeoclimate
data}\label{vegetation-and-palaeoclimate-data}

We sourced all community and climate records from \textbf{VegVault}
(version 1.0.0; {[}cite{]}), a harmonised SQLite database that
integrates fossil pollen archives with co-registered palaeoclimate
reconstructions and plant functional traits. Sourcing all three data
types required by this study --- vegetation community records,
co-registered palaeoclimate predictors, and plant functional traits ---
from a single harmonised repository offers two critical advantages.
First, it guarantees consistent site identifiers and temporal grids
across all inputs, which is a prerequisite for the aligned multi-species
modelling framework used here. Second, because VegVault is versioned and
publicly archived, the exact set of records used in this study can be
retrieved by any researcher with access to the same database version,
making the data inputs fully reproducible. Fossil pollen records are
drawn from the Neotoma Paleoecology
Database\textsuperscript{\textbf{Williams2018?}} and FOSSILPOL
{[}cite{]}; plant functional traits are drawn from
TRY\textsuperscript{\textbf{Kattge2020?}}. Palaeoclimate predictors are
CHELSA-TraCE21k bioclimatic
variables\textsuperscript{\textbf{Karger2021?}} co-registered in
VegVault at the site and time-step level.

We selected seven CHELSA bioclimatic variables as candidate abiotic
predictors: BIO1 (mean annual temperature), BIO4 (temperature
seasonality), BIO6 (minimum temperature of the coldest month), BIO12
(annual precipitation), BIO15 (precipitation seasonality), BIO18
(precipitation of the warmest quarter), and BIO19 (precipitation of the
coldest quarter). These variables represent the primary dimensions of
climate space known to govern plant distributions --- mean thermal
energy, climatic extremes, and water availability across seasons ---
while omitting correlated secondary variables to limit redundancy before
formal collinearity screening. Full descriptions and units of all
bioclimatic variables are provided in Supplementary Table S{[}TODO{]}.

The geographic and temporal coverage of extracted records (number of
fossil pollen sites per 500-year time slice and per continental region)
is summarised in Supplementary Figure S{[}TODO{]}.

\subsubsection{Taxonomic harmonisation}\label{taxonomic-harmonisation}

Fossil pollen records originate from multiple source databases that use
different original nomenclature systems and taxonomic conventions; a
unified classification is therefore a prerequisite for any
cross-database comparison. The guiding principle of the harmonisation
was to resolve each pollen morphotype to the \textbf{finest taxonomic
rank the morphological evidence supports}, rather than imposing a fixed
target rank across all morphotypes. For each morphotype, we first
attempted to assign a genus-level identification against the GBIF
taxonomic backbone {[}cite{]}, which provides a globally consistent,
community-maintained reference taxonomy. Where the morphotype was too
broadly defined morphologically to support a confident genus assignment
--- typically for compound types that aggregate morphologically similar
pollen from several genera --- the record was retained at family rank
rather than discarded, preserving ecological information that would
otherwise be lost entirely. We excluded non-plant records (e.g., fungal
spores and fern spores flagged in source databases). We supplemented the
backbone matching with a curated project-specific correction table that
records manually validated taxonomic decisions for morphotypes known to
be mis-classified or absent from the backbone, ensuring that recurring
edge cases are resolved consistently across all analyses. All
classifications were reviewed and validated by expert botanists to catch
systematic errors arising from regional nomenclatural traditions or
legacy morphotype names not covered by the backbone. Because the taxa
present in any given analysis depend on the geographic extent and
temporal window queried, we compiled a separate harmonisation table for
each continental region rather than a single global table. We provide
these tables as Supplementary Tables S{[}TODO{]} (one per continental
region).

\subsubsection{Multi-scale analytical
framework}\label{multi-scale-analytical-framework}

Co-occurrence patterns are shaped by processes operating at different
spatial extents, across geological time, and at different levels of
biological organisation. Because the relative importance of these
processes is itself expected to vary across scales, pooling all data at
a single arbitrary scale risks masking or confounding mechanisms that
are only visible at specific extents or resolutions. We therefore
applied the same modelling and analysis workflow at multiple levels
along three axes: \textbf{spatial} (continental, regional, and local
extents), \textbf{temporal} (40 successive 500-year slices from the Last
Glacial Maximum to the present), and \textbf{taxonomic} (genus, family,
and functional-type level). We first describe the shared pre-processing
and modelling steps common to all analyses, then detail how each axis
was defined and the reasoning behind specific choices.

\paragraph{Shared pre-processing}\label{shared-pre-processing}

\textbf{Community data preparation.} We applied several pre-processing
steps to the raw community matrices to prepare them for modelling and
ensure comparability across analyses. We converted raw pollen counts to
proportional abundance (relative pollen percentage). To standardise
temporal resolution across irregular sampling intervals (typical for
palaeoecological datasets), we then linearly interpolated records to a
common 500-year temporal grid. We considered binning but preferred
interpolation to preserve the full temporal resolution of the data and
avoid artefacts introduced by arbitrary bin edges. We excluded taxa
contributing less than 1 \% of the pollen sum within any individual
sample from that sample's record, because such occurrences are below the
detection limit reliably resolvable from pollen counts and contribute
predominantly noise to co-occurrence estimates. We used presence/absence
rather than continuous abundance as the response variable, because
pollen abundances are strongly influenced by taxon-specific pollen
production rates and dispersal distances --- biases that are effectively
removed by binarisation, leaving a response that more directly reflects
co-occurrence {[}TODO: cite{]}. Finally, to ensure sufficient data
density for reliable model estimates, we excluded spatial units with too
few sites or too few qualifying taxa after filtering. We calibrated the
exact thresholds separately for each spatial unit, because the number of
sites and taxa varies across units and the optimal thresholds depend on
the data structure of each unit.

\textbf{Abiotic predictor selection.} Collinear predictors inflate the
variance of coefficient estimates and may result in overfitting, because
variance attributable to one predictor is spuriously shared with its
collinear counterparts. We therefore screened candidate bioclimatic
variables for redundancy using a dual-filtering approach that combines
pairwise Pearson correlation and Variance Inflation Factor (VIF) {[}cite
Benito 2025{]}. We calibrated both thresholds automatically to the
correlation structure of each specific dataset rather than setting fixed
global values, because optimal thresholds depend on the number and
collinearity of predictors, which vary across spatial and temporal
units. When two predictors were redundant, we preferentially retained
the one with greater univariate predictive power, ensuring that the
non-redundant subset maximises ecological information.

\textbf{Variable scaling.} We centred and z-score scaled bioclimatic
predictors to a common variance, which is required for stable
gradient-descent optimisation and allows coefficient magnitudes to be
compared across predictors with different natural units. In analyses
with age as a predictor, we centred age but did not variance-scale it,
so that interaction coefficients retain their natural interpretation in
units of years. We projected site coordinates to a metric coordinate
system (ETRS89-LAEA, EPSG:3035) so that geographic distances are
measured in kilometres rather than decimal degrees.

\textbf{Moran Eigenvector Maps (MEMs).} Spatial autocorrelation in
species distributions arises from unmeasured environmental gradients,
dispersal limitation, and historical contingency --- processes that are
not captured by the bioclimatic predictors but would inflate the biotic
residual (B fraction) of the variance decomposition if left unaccounted
for. We represented large-scale spatial structure using Moran
Eigenvector Maps (MEMs), which are eigenvectors of a spatial weights
matrix and span the full range of positive spatial autocorrelation
present in the data {[}cite{]}. Including MEMs as predictor variables in
the spatial component of the model approximates a spatial random field
without the computational cost of full Gaussian process random effects,
which would be prohibitive for the large community matrices. For purely
spatial analyses, we derived MEMs from a two-dimensional geographic
distance matrix among sites. For spatiotemporal analyses, we used a
three-dimensional distance matrix (combining longitude, latitude, and
age) instead, so that the eigenvectors capture joint spatiotemporal
structure rather than spatial structure alone. We reduced the number of
retained eigenvectors at finer spatial scales, reflecting the lower
structural complexity of smaller analysis units.

\paragraph{Shared modelling workflow}\label{shared-modelling-workflow}

\textbf{Joint species distribution model.} We modelled plant
co-occurrence using a Joint Species Distribution Model
(jSDM)\textsuperscript{\textbf{Warton2015?}}, which fits all species
simultaneously within a single likelihood framework and estimates the
residual covariance among species after accounting for shared
environmental and spatial responses. We chose the jSDM framework over
single-species distribution models because it explicitly models the
residual inter-specific covariance --- that is, co-occurrences that
cannot be attributed to shared responses to environment or space ---
which is the central quantity of interest in this study. We performed
all analyses using the R package
\{sjSDM\}\textsuperscript{\textbf{Pichler2021?}}, which fits jSDMs via
stochastic gradient descent with an optional GPU backend, making it
feasible for the large community matrices (up to several hundred taxa ×
thousands of sites) analysed here.

Each taxon's occurrence was modelled as the linear combination of three
additive components:

\begin{itemize}
\tightlist
\item
  \textbf{Environment (A)} --- the non-collinear bioclimatic predictors,
  with age × climate interaction terms to allow the strength of climate
  filtering to vary across the Last Glacial--Interglacial cycle. only
  the interaction of age with climate predictors was included, the
  direct age term was excluded to avoid absorbing broad temporal trend
  into the environment component.
\item
  \textbf{Space (S)} --- MEM predictors included as fixed effects to
  absorb spatially (or spatiotemporally) structured variation not
  explained by the climate predictors, thereby preventing unmeasured
  spatial processes from inflating the biotic residual.
\item
  \textbf{Biotic residual (B)} --- a shared low-rank latent covariance
  matrix that captures the residual pairwise associations among taxa
  after accounting for their shared responses to environment and space.
  Positive residual covariance indicates that two taxa co-occur more
  often than expected given environment and spatial structure; negative
  covariance indicates mutual exclusion.
\end{itemize}

\textbf{Model fit and convergence.} Because the primary goal of this
study is variance decomposition rather than out-of-sample prediction, we
assessed model performance on the training data using Nagelkerke R² and
per-species AUC, accuracy, and log-loss. Convergence was confirmed by
inspecting the training loss trajectory: we considered a model converged
when the linear trend of the loss curve was negligible and successive
steps produced only marginal changes, indicating that further
optimisation would not meaningfully alter the parameter estimates.

\textbf{Variance decomposition.} To attribute the variance in community
composition to ecologically interpretable sources, we applied a Monte
Carlo ANOVA decomposition to each fitted jSDM. The approach repeatedly
draws from the fitted model with one or more components omitted and
quantifies how much variation is lost, yielding seven non-overlapping
Venn-diagram fractions:

\begin{itemize}
\tightlist
\item
  \textbf{A} --- variance uniquely attributable to the abiotic
  (environment) component;
\item
  \textbf{S} --- variance attributable to the spatial component (MEMs);
\item
  \begin{itemize}
  \tightlist
  \item
    \textbf{B} --- variance attributable to the residual biotic
    covariance (co-occurrence patterns unexplained by environment or
    space, interpreted as a signal of potential biotic
    interactions\textsuperscript{\textbf{Warton2015?}});
  \end{itemize}
\item
  \textbf{AB, AS, BS, ABS} --- variance jointly attributable to pairs or
  all three components.
\end{itemize}

We set fractions that were negative --- an artefact of Monte Carlo
approximation error --- to zero, as negative explained variance has no
ecological interpretation.

\paragraph{Spatial axis}\label{spatial-axis}

To isolate the effect of geographic extent on co-occurrence patterns, we
applied the shared workflow independently at three nested spatial scales
while keeping the temporal and taxonomic dimensions constant. The three
scales were chosen to span the range from macroecological processes ---
where broad climatic gradients dominate co-occurrence --- to local
processes, where dispersal limitation and microhabitat filtering are
more prominent:

\begin{itemize}
\tightlist
\item
  \textbf{Continental} --- three units: Europe (−10° to 40°E, 35° to
  70°N), North America (−130° to −60°W, 30° to 70°N), and northern Asia
  (60° to 140°E, 50° to 75°N).
\item
  \textbf{Regional} --- approximately 30 geographic units per continent,
  each spanning roughly 20° longitude × 10°--15° latitude.
\item
  \textbf{Local} --- approximately 100 fine-grain units per continent,
  each spanning roughly 5° longitude × 5° latitude.
\end{itemize}

We compared variance fractions across scales to identify which
components of co-occurrence signal are scale-invariant and which are
scale-dependent.

\textsubscript{Source:
\href{https://OndrejMottl.github.io/BIODYNAMICS-vegetation-cooccurrence/manuscript.qmd.html}{Article
Notebook}}

Three scale-specific thresholds governed which spatial units were
retained for model fitting. We excluded a taxon from a spatial unit if
it appeared in fewer than 5 distinct sediment cores. We further excluded
taxa if they appeared in fewer than 10 distinct site--age combinations
at the continental and regional scales, or fewer than 5 at the local
scale, reflecting the reduced number of sites per unit at finer grains.
We excluded a spatial unit entirely if fewer than 5 qualifying taxa
remained after all pre-processing at the continental scale; the
corresponding lower bounds were 4 taxa at the regional scale and 3 taxa
at the local scale.

\paragraph{Temporal axis}\label{temporal-axis}

To test whether the structure of plant co-occurrence has changed through
time --- and in particular whether the large climate changes since the
Last Glacial Maximum are associated with a reorganisation of
co-occurrence networks --- we applied the shared workflow independently
to each 500-year time slice across the full 20,000-year window. We
fitted a separate jSDM to each time slice rather than a single pooled
model with a temporal random effect. Prior to model fitting, we computed
co-occurrence network metrics directly from the binarised community
matrix at each time slice to characterise the structural properties of
the co-occurrence network independently of any model assumptions. We
then aggregated results from successive time slices into temporal
trajectories of variance fractions and network metrics, revealing how
the relative importance of environment, space, and biotic residuals ---
and the structure of the co-occurrence network itself --- has evolved
over geological time. We ran analyses for three continental regions:
Europe, North America, and northern Asia.

\textsubscript{Source:
\href{https://OndrejMottl.github.io/BIODYNAMICS-vegetation-cooccurrence/manuscript.qmd.html}{Article
Notebook}}

We applied the same three data-quality thresholds at each time slice
within each continental region. We excluded a taxon from a given time
slice if it appeared in fewer than 5 distinct sediment cores, or in
fewer than 10 distinct site--age combinations. We excluded a time slice
entirely if fewer than 5 qualifying taxa remained after all
pre-processing.

\paragraph{Taxonomic axis}\label{taxonomic-axis}

The level at which taxa are defined determines which ecological
processes are legible in co-occurrence patterns: genus-level
co-occurrence may reflect fine-grained competitive or facilitative
interactions, whereas broader groupings may reveal assembly rules driven
by shared functional strategies. To test how strongly co-occurrence
patterns depend on the level of biological organisation, we collapsed
the community matrices to three taxonomic resolutions: \textbf{Genus},
\textbf{Family} and \textbf{Functional type (FT)} --- ecologically
coherent groups derived from functional trait clustering (see
\emph{Functional-type classification} below).

We applied the same jSDM pipeline at each resolution without
modification, enabling direct comparison of variance fractions across
taxonomic scales. We evaluated the taxonomic axis only at the spatial
scale, fitting each resolution independently per spatial unit.

\subsubsection{Functional-type
classification}\label{functional-type-classification}

\textbf{Trait data and quality control.} Plant functional traits
characterise the ecological strategies of taxa and provide a
biologically grounded basis for grouping species into functional types
that reflect shared responses to the environment rather than arbitrary
taxonomic boundaries. We extracted trait data from VegVault for all
available trait domains (see Supplementary Table S{[}TODO{]} for the
full list of traits and units). Because trait databases aggregate
measurements from many sources with varying units and methodologies, raw
values frequently contain outliers and unit errors that would distort
the distance matrix used for clustering. We therefore applied a
two-stage quality-control workflow to the trait data, reviewing and
correcting all flagged records manually. In the first stage, we
identified outlier taxa using an inter-quartile range criterion and
corrected or excluded them after review. In the second stage, we applied
a further quality-control pass after taxonomic harmonisation to remove
trait values assigned to ambiguous or unresolvable morphotypes. We
aggregated validated values to taxon-level medians per trait domain and
assembled them into a taxon × trait matrix for clustering.

\textbf{Distance computation and clustering.} We derived functional
types separately for each continental unit (Europe, North America,
northern Asia) to ensure that the groupings reflect the functional
diversity actually present in each region rather than a globally imposed
classification. We quantified pairwise dissimilarities among taxa using
the Gower distance metric, which is appropriate for plant trait data
because it handles mixed variable types (continuous, binary, ordinal)
and missing values natively, without requiring imputation or prior
standardisation of trait distributions. We then grouped taxa by Ward's
minimum-variance hierarchical clustering applied to the Gower distance
matrix. We chose Ward's method because it minimises within-cluster total
variance at each merge step, producing compact, well-separated clusters
that are more likely to correspond to ecologically coherent functional
strategies than linkage methods that optimise only local cohesion.

\textbf{Selection of the number of functional types.} We determined the
number of functional types \emph{k} by a dual criterion designed to
balance cluster quality with data availability for downstream modelling.
For each candidate \emph{k} in the range \emph{k}\textsubscript{min} (=
5) to \emph{k}\textsubscript{max} (= 30), we passed the resulting
FT-level community matrix through the same data-filtering steps
(rare-type removal, minimum-site and minimum-sample filters, and
binarisation). We considered viable only \emph{k} values that left at
least a minimum number of non-constant FT columns after filtering,
because a \emph{k} that produces many rare or invariant functional types
provides insufficient data for stable model estimation regardless of its
clustering quality. Among the viable candidates, we selected the
\emph{k} with the highest mean silhouette width --- a model-free measure
of cluster cohesion and separation that does not depend on the
clustering algorithm's internal objective --- as the final number of
functional types. This two-stage criterion prevents the selection
algorithm from choosing a fine-grained \emph{k} that scores well on
clustering metrics but collapses to an inadequate number of informative
functional types after real-data filtering. We then assigned each taxon
its functional-type label and collapsed the community matrix to FT level
by taking the union of presences across all member taxa within each site
and time point.

\subsubsection{Computational
reproducibility}\label{computational-reproducibility}

Reproducibility of the full analysis was a primary design goal, because
the study spans hundreds of independent models across three continents,
three spatial scales, and 40 time slices --- a scope at which silent
failures and undocumented parameter choices are difficult to detect
after the fact. We implemented the entire analytical workflow --- from
data extraction through modelling, variance decomposition, and figure
generation --- as a declarative, dependency-tracked pipeline using the
\{targets\} R package\textsuperscript{\textbf{Landau2021?}}, which
caches the output of every computation and automatically re-runs only
those steps that are invalidated by an upstream change in code or data.
This design makes every result traceable to the exact code and inputs
that produced it and ensures that partial re-runs after code edits
produce the same output as a full run from scratch. We locked package
versions using \{renv\}\textsuperscript{\textbf{Ushey2023?}}, so that
the full computational environment --- including all R package versions
and their dependencies --- can be recreated identically on any machine.

\textsubscript{Source:
\href{https://OndrejMottl.github.io/BIODYNAMICS-vegetation-cooccurrence/manuscript.qmd.html}{Article
Notebook}}

The analytical codebase comprises 113 modular R functions, each in its
own file, spanning data extraction, taxonomic harmonisation, community
pre-processing, collinearity screening, spatial eigenvector computation,
model fitting, variance decomposition, functional-type classification,
and visualisation. Each function has a dedicated unit-test file (112
test files in total), and the test suite achieves 92 \% line coverage
across all function files as measured by \{covr\} {[}TODO: cite{]}. All
113 functions are documented and are publicly accessible through a
dedicated project documentation website {[}GitHub Pages URL; TODO{]},
which provides searchable reference pages for every function alongside
narrative descriptions of each analysis phase.

All analysis code, pipeline definitions, configuration files, and
instructions for reproducing the environment are publicly available at
{[}GitHub URL / Zenodo DOI; TODO{]}.

\subsection{Results}\label{results}

\phantomsection\label{cell-fig-cooccurrence-overview}
\textsubscript{Source:
\href{https://OndrejMottl.github.io/BIODYNAMICS-vegetation-cooccurrence/manuscript.qmd.html}{Article
Notebook}}

\subsubsection{Data overview}\label{data-overview}

TODO: Report number of sites, taxa, time windows, and samples retained
after filtering (use inline \texttt{r} expressions --- no hardcoded
numbers).

\subsubsection{Model performance}\label{model-performance}

TODO: Report predictive performance metrics (AUC, RMSE, etc.).

\subsubsection{Co-occurrence patterns}\label{co-occurrence-patterns}

TODO: Describe spatial and temporal patterns in species associations.

\subsubsection{Environmental drivers}\label{environmental-drivers}

TODO: Describe ANOVA variance partitioning results.

\subsection{Discussion}\label{discussion}

TODO: Write discussion.

\subsection{Acknowledgements}\label{acknowledgements}

TODO: Add acknowledgements. Funding sources, data providers, computing
resources, field/lab assistance.

\subsection{Author Contributions}\label{author-contributions}

\begin{longtable}[]{@{}ll@{}}
\toprule\noalign{}
Role & Contributors \\
\midrule\noalign{}
\endhead
\bottomrule\noalign{}
\endlastfoot
Conceptualisation & \\
Data curation & \\
Formal analysis & \\
Funding acquisition & \\
Investigation & \\
Methodology & \\
Project administration & \\
Resources & \\
Software & \\
Supervision & \\
Validation & \\
Visualisation & \\
Writing -- original draft & \\
Writing -- review \& editing & \\
\end{longtable}

\subsection{Data Availability}\label{data-availability}

TODO: Describe data availability. Reference VegVault, pipeline outputs,
and any archived datasets.

\subsection{Code Availability}\label{code-availability}

The analysis code is available at
\url{https://github.com/OndrejMottl/BIODYNAMICS-vegetation-cooccurrence}
and archived at TODO (Zenodo DOI to be added upon acceptance).

\subsection{References}\label{references}

\phantomsection\label{refs}

\section*{Supplementary Information}\label{supplementary-information}
\addcontentsline{toc}{section}{Supplementary Information}

\subsection{Supplementary Figures}\label{supplementary-figures}

\subsubsection*{Figure S1 --- TODO title}\label{figure-s1-todo-title}
\addcontentsline{toc}{subsubsection}{Figure S1 --- TODO title}

\phantomsection\label{cell-fig-s1}
\textsubscript{Source:
\href{https://OndrejMottl.github.io/BIODYNAMICS-vegetation-cooccurrence/manuscript.qmd.html}{Article
Notebook}}

\subsection{Supplementary Tables}\label{supplementary-tables}

\subsubsection*{Table S1 --- TODO title}\label{table-s1-todo-title}
\addcontentsline{toc}{subsubsection}{Table S1 --- TODO title}

\subsection{Supplementary Methods}\label{supplementary-methods}

TODO: Add any extended methodological details not included in the main
Methods section (e.g., sensitivity analyses, model diagnostics,
additional filtering criteria).




\end{document}
